\par\medskip
\noindent
5.3. Utilizando la correspondencia de ortogonalidad entre los
prosubgrupos estrictos de Barsotti--Tate (en el sentido precisado más
arriba) de los progrupos duales $T_\ell(A_K)$ y $T_\ell(A'_K)$, se
deduce formalmente de (5.2.2) que el emparejamiento (1.0.2) entre estos
progrupos induce un emparejamiento
\begin{equation}\tag{5.3.2}
 T_\ell(A_K)^{ab}\times T_\ell(A'_K)^{ab}
 \longrightarrow T_\ell(\mathbb G_{mK})
\end{equation}
que es una \underline{dualidad perfecta}, es decir, identifica cada uno
de los dos factores de la izquierda con el dual del otro (en el sentido
de 5.2.3).\footnote{Los superíndices $ab$ denotan las \emph{partes
abelianas}, el cociente de la parte fija por la parte tórica.} Resulta
también claro que (5.3.2) está inducido por el emparejamiento
\begin{equation}\tag{5.3.3}
 T_\ell(A^o)^{ab}\times T_\ell(A^{\prime o})^{ab}
 \longrightarrow T_\ell(\mathbb G_{mS})
\end{equation}
obtenido del emparejamiento (5.2.4.2) pasando a los cocientes. Puesto
que esta es una dualidad perfecta sobre las fibras genéricas, es decir,
puesto que el homomorfismo correspondiente
\begin{equation}\tag{5.3.4}
 T_\ell(A^o)^{ab}
 \longrightarrow D\bigl(T_\ell(A^{\prime o})^{ab}\bigr)
\end{equation}
es un isomorfismo sobre las fibras genéricas, es natural querer deducir
que lo es sobre todo $S$, o equivalentemente sobre las fibras
especiales. Esta deducción es trivialmente válida si $\ell\neq p$.
También lo es, al menos en características desiguales cuando $\ell=p$,
por un profundo teorema de Tate [33, Th. 4]. No obstante, la
estableceremos directamente más adelante, en 7.4.6, sin utilizar el
teorema de Tate.

Para precisar el significado geométrico del resultado que queremos
demostrar, obsérvese que el emparejamiento obtenido de (5.3.3) al pasar
a las fibras especiales es precisamente el emparejamiento
\begin{equation}\tag{5.3.5}
 T_\ell(B_o)\times T_\ell(B'_o)
 \longrightarrow T_\ell(\mathbb G_{mk})
\end{equation}
deducido de la biextensión $W_o$ de $(A_o^o,A_o^{\prime o})$ pasando
al cociente, donde $B_o$ y $B'_o$ denotan respectivamente las partes
abelianas de $A_o^o$ y $A_o^{\prime o}$. Por VIII 4.8, la biextensión
$W_o$ procede de una biextensión $W_o^{ab}$ de $(B_o,B'_o)$ por
$\mathbb G_{mk}$, única salvo isomorfismo único. Así, el emparejamiento
(5.3.5) puede interpretarse también como el deducido de la
correspondencia divisorial $W_o^{ab}$, y el homomorfismo correspondiente
\begin{equation}\tag{5.3.6}
 T_\ell(B_o)
 \longrightarrow D\bigl(T_\ell(B'_o)\bigr)
 \simeq T_\ell\bigl(D(B'_o)\bigr)
\end{equation}
(donde, en el último término, $D$ denota el esquema abeliano dual) es
precisamente el morfismo obtenido aplicando funtorialmente $T_\ell$ al
homomorfismo
\begin{equation}\tag{5.3.7}
 B_o\longrightarrow D(B'_o)
\end{equation}
definido por $W_o^{ab}$. En consecuencia, según el diccionario usual
(cf. [5]), decir que (5.3.5) es una dualidad perfecta equivale también
a decir que el núcleo de (5.3.7), que es un grupo finito sobre $k$ por
1.5, posee componente $\ell$-primaria nula. Decir que esto se cumple
para todo $\ell$ significa, por tanto, que la isogenia (5.3.7) es un
isomorfismo, es decir, que la correspondencia divisorial $W_o^{ab}$ es
una dualidad perfecta entre $B_o$ y $B'_o$. Las consideraciones
precedentes solo muestran que el núcleo de (5.3.7) es un $p$-grupo.
Enunciamos, pues, el resultado más completo.

\medskip
\noindent
\underline{Teorema} 5.4. \underline{Sea $S$ un rasgo, y sea $A_K$ un
esquema abeliano sobre \sisi{$S$}{$K$} que posee reducción semiestable
sobre $S$. Sean $B_o$ y $B'_o$ las partes abelianas de los modelos de
Néron de $A_K$ y del esquema abeliano dual $A'_K$, respectivamente.
Entonces la correspondencia divisorial canónica $W_o^{ab}$ sobre
$(B_o,B'_o)$, deducida por VIII 4.8 de la biextensión canónica $W_o$ de
$(A_o^o,A_o^{\prime o})$ (la fibra especial de la biextensión $W$ de
1.4), es una dualidad perfecta entre $B_o$ y $B'_o$; es decir, define
un isomorfismo de cada uno de estos esquemas abelianos con el dual del
otro.}

La demostración de 5.4 en el caso general se dará en la sección 7, como
consecuencia de una construcción muy natural e importante debida a
Raynaud. Limitémonos a observar aquí que, para demostrar 5.4, puede
suponerse que $S$ es completo, por 3.3.
